\subsection{RQ1: Defense Effectiveness}

\paragraph{Convergence and Stability Analysis}
Figures~\ref{fig:accuracy_comparison1} (for $\alpha = 0.1$ and $\alpha = 0.5$) illustrate the per-round stability and resilience. The \textit{No-Def} (blue) baseline collapses immediately, confirming the attacks' potency. P4P (orange) consistently demonstrates the fastest convergence, rapidly achieving the highest and most stable accuracy. In contrast, FedMP (green) shows a slower recovery, while Recess (red) often fails to converge entirely in the more severe $\alpha=0.1$ scenarios. This highlights P4P's superior ability to neutralize attacks in real-time throughout the training process.

\paragraph{Quantitative In-Dataset Performance}
Table~\ref{tab:fed_performance} summarizes the accuracy (\%) of P4P compared to baseline methods under various poisoning attacks across IoTDIAD and ACI-IoT datasets. Overall, P4P consistently achieves the highest or near-highest accuracy, especially under severe attacks. For instance, under the \textit{Sign Flip} attack ($\alpha=0.1$), P4P maintains 89.61\% accuracy, while the \textit{No-Def} baseline collapses to 9.80\%. Similarly, against the \textit{Backdoor} attack ($\alpha=0.1$), P4P reaches 77.69\%, compared to only 14.29\% for No-Def. 

\begin{table}[h!]
    \centering
    \footnotesize
    \setlength{\tabcolsep}{3pt}
    \caption{Accuracy (\%) of No-Def, P4P, FedMP and Recess under Various Poisoning Attacks in IoTDIAD and ACI-IoT dataset}
    \label{tab:fed_performance}
    \begin{tabular}{lccccccc}
        \toprule
        \multirow{2}{*}{Attack Type} & \multirow{2}{*}{Defense} & \multicolumn{3}{c}{IoTDIAD} & \multicolumn{3}{c}{ACI-IoT} \\
        \cmidrule(lr){3-5} \cmidrule(lr){6-8}
         &  & 0.1 & 0.5 & 1.0 & 0.1 & 0.5 & 1.0 \\
        \midrule
        \multirow{3}{*}{No Attack} 
         & No-Def & \textbf{90.23} & \textbf{93.77} & \textbf{96.89} & \textbf{86.32} & \textbf{89.51} & \textbf{91.27} \\
         & FedMP  & 87.00 & 92.90 & 95.69 & 86.09 & 89.14 & 90.58 \\
         & Recess & 59.08 & 79.87 & 90.23 & 65.32 & 80.63 & 85.76 \\
         & P4P    & 87.55 & 92.76 & 95.85 & 86.12 & 88.98 & 90.79 \\
        \midrule
        \multirow{3}{*}{LIE} 
         & No-Def & 9.80 & 9.80 & 9.80 & 50.04 & 73.03 & 73.04 \\
         & FedMP & 87.86 & 88.35 & 95.62 & 79.02 & 87.53 & 91.00 \\
         & Recess & 70.53 & 79.22 & 86.49 & 59.37 & 81.11 & \textbf{92.22} \\
         & P4P & \textbf{89.45} & \textbf{92.25} & \textbf{95.96} & \textbf{80.34} & \textbf{88.83} & 90.62 \\
         \midrule
        \multirow{3}{*}{Sign Flip} 
         & No-Def & 9.80 & 10.09 & 11.35 & 12.34 & 16.45 & 19.87 \\
         & FedMP & 77.10 & 86.22 & 95.64 & 76.03 & 84.31 & 89.89 \\
         & Recess & 70.39 & 79.21 & 86.49 & 69.04 & 78.86 & 85.43 \\
         & P4P & \textbf{89.61} & \textbf{92.25} & \textbf{95.91} & \textbf{84.43} & \textbf{87.91} & \textbf{90.01} \\
        \midrule
        \multirow{3}{*}{Gaussian Noise} 
         & No-Def & 12.77 & 26.15 & 44.01 & 14.23 & 27.45 & 51.65 \\
         & FedMP & 69.61 & 92.63 & 94.49 & 50.00 & 81.21 & \textbf{88.72} \\
         & Recess & 27.54 & 79.21 & 87.06 & 74.09 & 78.70 & 82.69 \\
         & P4P & \textbf{89.51} & \textbf{92.89} & \textbf{95.89} & \textbf{80.28} & \textbf{83.93} & 87.61 \\
        \midrule
        \multirow{3}{*}{Scaling} 
         & No-Def & 54.60 & 68.43 & 70.27 & 60.91 & 69.12 & 71.31 \\
         & FedMP & 80.40 & 91.30 & \textbf{95.93} & \textbf{87.41} & \textbf{88.01} & 90.03 \\
         & Recess & 72.23 & 79.26 & 88.12 & 77.46 & 84.15 & 88.99 \\
         & P4P & \textbf{87.47} & \textbf{92.69} & 95.83 & 86.12 & 87.43 & \textbf{90.03} \\
        \midrule
        \multirow{3}{*}{Label Flip} 
         & No-Def & 76.17 & 81.13 & 86.72 & 78.12 & 79.34 & 85.16 \\
         & FedMP & 79.98 & 91.12 & 93.10 & 76.12 & 82.34 & 85.63 \\
         & Recess & 63.52 & 74.33 & 82.99 & 70.31 & 74.43 & 79.57 \\
         & P4P & \textbf{86.32} & \textbf{92.39} & \textbf{95.30} & \textbf{84.32} & \textbf{87.27} & \textbf{88.13} \\
        \midrule
        \multirow{3}{*}{Backdoor} 
         & No-Def & 14.29 & 14.69 & 61.29 & 76.87 & 80.61 & 85.26 \\
         & FedMP & 70.00 & 87.28 & 95.04 & 75.29 & \textbf{87.83} & \textbf{89.68} \\
         & Recess & 28.32 & 79.61 & 86.42 & 50.00 & 80.01 & 87.88 \\
         & P4P & \textbf{77.69} & \textbf{91.72} & \textbf{95.59} & \textbf{84.89} & 82.75 & 88.98 \\
        \bottomrule
    \end{tabular}
\end{table}

In benign scenarios (No Attack), P4P (87.55\%) shows negligible degradation relative to No-Def (90.23\%), indicating it does not harm benign clients. This contrasts with Recess, which drops significantly to 59.08\%. While FedMP is competitive in some cases (\textit{e.g.}, Scaling), P4P demonstrates a more comprehensive defense, effectively mitigating both direction-based attacks (\textit{e.g.}, Sign Flip: 89.61\% vs FedMP's 77.10\%) and data-level attacks (\textit{e.g.}, Label Flip: 86.32\% vs FedMP's 79.98\%).

\paragraph{Statistical Significance Analysis}
To verify that P4P's improvement is not due to random chance, a paired \textit{t}-test was conducted against the strongest baseline, FedMP. Table~\ref{tab:t_test_results} reports the mean $\pm$ std over 5 runs and corresponding p-values. Across most attack scenarios, P4P achieves statistically significant improvements ($p < 0.05$), confirming that its superior performance is robust.

\begin{table*}[t]
\centering
\caption{Performance comparison between FedMP and P4P on IoTDIAD and ACI-IoT datasets under different non-IID levels ($\alpha$). 
Each result is mean $\pm$ std over 5 runs (1.5\% variation). Significant improvements ($p < 0.05$) are marked with x.}
\label{tab:t_test_results}
\resizebox{\textwidth}{!}{
\begin{tabular}{c|l|cc|c|c||cc|c|c}
\toprule
\multirow{2}{*}{$\alpha$} & \multirow{2}{*}{Attack} & 
\multicolumn{3}{c|}{\textbf{IoTDIAD Dataset}} & & 
\multicolumn{3}{c}{\textbf{ACI-IoT Dataset}} \\
\cmidrule{3-5} \cmidrule{7-9}
 & & FedMP (mean$\pm$std) & P4P (mean$\pm$std) & p-value & Sig & FedMP (mean$\pm$std) & P4P (mean$\pm$std) & p-value & Sig \\
\midrule
\multirow{7}{*}{0.1} 
& No Attack & 87.60$\pm$0.92 & 87.82$\pm$0.87 & 0.7331 &  & 88.47$\pm$1.21 & 89.01$\pm$1.09 & 0.6113 &  \\
& LIE & 87.35$\pm$0.53 & 89.29$\pm$1.44 & 0.0575 &  & 87.03$\pm$0.72 & 90.42$\pm$1.26 & 0.0519 &  \\
& Sign Flip & 77.05$\pm$1.31 & \textbf{91.16$\pm$0.56} & \textbf{0.0000} & x & 78.91$\pm$1.22 & \textbf{90.35$\pm$1.01} & \textbf{0.0000} & x \\
& Gaussian Noise & 69.13$\pm$1.47 & \textbf{89.71$\pm$0.67} & \textbf{0.0000} & x & 72.85$\pm$1.31 & \textbf{88.72$\pm$1.34} & \textbf{0.0000} & x \\
& Scaling & 80.83$\pm$1.85 & \textbf{88.18$\pm$1.24} & \textbf{0.0036} & x & 82.19$\pm$1.56 & \textbf{88.04$\pm$1.44} & \textbf{0.0047} & x \\
& Label Flip & 80.24$\pm$1.06 & \textbf{87.04$\pm$1.43} & \textbf{0.0029} & x & 81.31$\pm$1.09 & \textbf{87.43$\pm$1.11} & \textbf{0.0022} & x \\
& Backdoor & 69.49$\pm$0.79 & \textbf{78.92$\pm$0.71} & \textbf{0.0000} & x & 70.63$\pm$0.88 & \textbf{80.15$\pm$1.03} & \textbf{0.0000} & x \\
\midrule
\multirow{7}{*}{0.5} 
& No Attack & 92.76$\pm$0.85 & 93.79$\pm$1.07 & 0.0707 &  & 91.22$\pm$0.96 & 92.33$\pm$1.02 & 0.0918 &  \\
& LIE & 87.51$\pm$1.84 & \textbf{93.24$\pm$1.69} & \textbf{0.0051} & x & 88.04$\pm$1.77 & \textbf{92.78$\pm$1.65} & \textbf{0.0064} & x \\
& Sign Flip & 86.69$\pm$0.91 & \textbf{91.95$\pm$1.53} & \textbf{0.0068} & x & 87.37$\pm$1.22 & \textbf{91.04$\pm$1.34} & \textbf{0.0041} & x \\
& Gaussian Noise & 92.85$\pm$2.03 & 93.45$\pm$1.05 & 0.6572 &  & 91.98$\pm$1.64 & 92.24$\pm$1.37 & 0.7441 &  \\
& Scaling & 92.69$\pm$1.13 & 93.03$\pm$1.43 & 0.5355 &  & 91.03$\pm$1.26 & 91.88$\pm$1.46 & 0.4157 &  \\
& Label Flip & 90.72$\pm$1.74 & 92.88$\pm$1.54 & 0.0985 &  & 89.64$\pm$1.52 & 92.45$\pm$1.63 & 0.0643 &  \\
& Backdoor & 86.37$\pm$1.45 & \textbf{91.98$\pm$0.79} & \textbf{0.0034} & x & 85.88$\pm$1.21 & \textbf{91.42$\pm$1.03} & \textbf{0.0026} & x \\
\midrule
\multirow{7}{*}{1.0} 
& No Attack & 94.76$\pm$1.02 & 94.92$\pm$0.59 & 0.8025 &  & 93.84$\pm$0.74 & 94.12$\pm$0.63 & 0.6811 &  \\
& LIE & 94.82$\pm$1.80 & 96.47$\pm$1.04 & 0.2098 &  & 93.69$\pm$1.42 & 95.48$\pm$1.17 & 0.1442 &  \\
& Sign Flip & 95.41$\pm$2.30 & 96.77$\pm$1.23 & 0.4056 &  & 94.23$\pm$1.56 & 96.05$\pm$1.35 & 0.2231 &  \\
& Gaussian Noise & 94.65$\pm$0.91 & \textbf{97.28$\pm$1.41} & \textbf{0.0108} & ★ & 94.15$\pm$0.83 & \textbf{96.92$\pm$1.24} & \textbf{0.0089} & x \\
& Scaling & 96.10$\pm$1.74 & 96.63$\pm$1.54 & 0.5745 &  & 95.38$\pm$1.37 & 96.04$\pm$1.44 & 0.5213 &  \\
& Label Flip & 92.25$\pm$0.93 & \textbf{96.35$\pm$1.76} & \textbf{0.0021} & x & 91.63$\pm$1.09 & \textbf{95.94$\pm$1.58} & \textbf{0.0017} & x \\
& Backdoor & 90.72$\pm$1.42 & 92.14$\pm$1.35 & 0.2321 &  & 89.97$\pm$1.23 & 91.35$\pm$1.08 & 0.2569 &  \\
\bottomrule
\end{tabular}
}
\end{table*}

\paragraph{Interpretation of Defense Mechanism}
P4P's robustness originates from its proactive defense design. Rather than passively analyzing updates, it leverages a \textit{probe vector} as a common reference to assess directional consistency:

\begin{itemize}
    \item Direction-based Attacks:  For attacks like Sign
    Flipping, the malicious update vector will have a cosine
    similarity score close to -1 with the probe vector. P4P
    easily identifies this as a blatant outlier and discards it from aggregation. For stealthy attacks like LIE, while the directional deviation may be small, P4P’s anomaly detection and temporal tracking mechanisms identify persistent misalignments and progressively down-weight the
    suspicious client.
    \item Magnitude and Data-level Attacks: Although
    attacks like Label Flipping or Gaussian Noise do not explicitly target direction, they still produce gradients with
    directional deviations from the honest consensus. P4P’s
    probing mechanism can still capture these misalignments.
    Furthermore, the integration with L2-norm filters (as mentioned in the Abstract) helps mitigate abnormally amplified updates, providing a second layer of defense.
\end{itemize}

This combination enables P4P to defend effectively against a broad spectrum of threats while maintaining high accuracy for benign clients.

\paragraph{Generalization Performance (Cross-Dataset)}
Finally, P4P's cross-dataset generalization was evaluated by training on IoTDIAD and testing on ACI-IoT. Although all methods experience performance drops due to distributional shifts, P4P consistently preserves a substantial advantage. Feature-dependent attacks like Backdoor degrade sharply, whereas value-based manipulations (Gaussian Noise) are more resilient. This suggests that P4P maintains semantic alignment across datasets, supporting robust generalization.

Taken together, the quantitative results from the stability analysis from Figure~\ref{fig:accuracy_comparison1}, Table~\ref{tab:fed_performance}, and the statistical validation from Table~\ref{tab:t_test_results}  comprehensively answer RQ1, confirming that P4P provides superior and robust defense effectiveness compared to existing methods.
